\section{Transfer to SGLang}
\label{app:e4}
The engine study evaluates the shared policy and accounting contract on a second serving
stack. It reuses H, H+guard, two deployment checkpoints, and two screening checkpoints
(Table~\ref{tab:e4identities}), with no retraining. Among the screening study's historically
accepted and rejected candidates under Lookup-69, the lowest mean physical-cost
candidate in each group supplies the two representatives. This fixed,
nonrandom cohort tests transfer of existing policies and predictions.

\begin{table}[t]
\centering
\small
\caption{Frozen PPO identities reused in the engine study. PPO-1/PPO-2 come from
the deployment study and PPO-7/PPO-9 from checkpoint screening. Rule parameters and guard safety estimates are unchanged.}
\label{tab:e4identities}
\setlength{\tabcolsep}{3pt}
\par\vspace{4pt}
\begin{tabular*}{\linewidth}{@{\extracolsep{\fill}}lll@{}}
\toprule
Label & Training environment & Seed / checkpoint\\
\midrule
PPO-1 & Lookup-69 & 101 / 40\\
PPO-2 & Lookup-69 & 103 / 100\\
PPO-7 & Lookup-30 & 105 / 35\\
PPO-9 & Lookup-30 & 102 / 80\\
\bottomrule
\end{tabular*}
\end{table}

\paragraph{Backend and execution.}
The deployment uses Qwen2.5-7B-Instruct on two A100 80GB PCIe GPUs,
one independent TP1/PP1 SGLang 0.4.8 instance per GPU, FP16, model length
4096, static-memory fraction 0.8, chunked-prefill size and maximum prefill
tokens 16,384, maximum running requests 1,024, and four local slots.
CUDA graphs and radix caching are disabled. SGLang uses PyTorch 2.7.1 and
Transformers 4.52.3; the CPU policy controller retains PyTorch 2.7.0.
Native token-ID input avoids retokenizing the frozen workloads. A recorded
compatibility fix supplies model EOS metadata when the tokenizer is not
initialized, preserving minimum-output-token masking. The fix, input-token
identities, output-budget checks, readiness probes, and verified process
release are bound in the source and execution manifests.

Recovery and steady reuse deployment workloads: respectively 358 and 944
online arrivals, plus 120 offline arrivals per window. Each of six arms
runs twice per window, with one all2 anchor per window: 24 policy runs
and two anchors. The observation interval remains 2100\,s, with the same
deadline, routing, guard, and scenario-price definitions. Separate runtime
identities are retained; the study is a transfer of the full serving stack,
not an isolated engine-performance comparison.

\paragraph{Startup measurement and frozen predictions.}
Twenty calibration cycles launch the second instance while the first
serves four concurrent 256-input/256-output-token background requests.
Timing runs from launch request to parent-verified readiness. All 20 cycles
complete, with median 27.5675401645\,s and range 26.713--29.165\,s;
setup timing is recorded separately. Lookup-SG substitutes that unrounded
median for Lookup-69's 69.219185665\,s startup proxy. Lookup-30 uses 30\,s.
All three share the inherited service lookup, 7.059\,s shutdown proxy,
policy identities, and accounting rules; the guard's safety estimates
remain unchanged. This calibrates startup, not the full engine service model.

Design and compatibility checks precede calibration; the GPU-side design
check is local. The 42 predictions (three models, seven policies including
all2, two windows) and measured startup parameter are frozen and independently
accepted on CPU before the formal runs. Physical repeats reuse the same
prediction, so the 24 comparisons per model represent 12 unique policy/window
predictions rather than 24 independent samples.

\paragraph{Results.}
All 26 final runs are technically valid. Both all2 anchors meet every deadline,
occupy 4200 GPU seconds, and cost 4.2000 scenario USD each. All online
requests are timely. H leaves one offline job unfinished in each steady
repeat; all other policy runs complete all work on time.
Tables~\ref{tab:e4costs} and~\ref{tab:e4resources} retain both repeats of
every policy/window pair. P1+ failures remain in the agreement and cost-error
denominators. Mean absolute error uses total scenario cost, including the
offline-miss penalty, consistently with the main paper.

\begin{table*}[t]
\centering
\small
\caption{Engine study predicted and observed total window costs in scenario USD.
R/S denote recovery/steady. A dagger marks predicted or observed P1+ failure;
unmarked entries satisfy P1+. Both all2 anchors cost 4.2000 and satisfy P1+.}
\label{tab:e4costs}
\setlength{\tabcolsep}{4pt}
\par\vspace{4pt}
\begin{tabular*}{0.96\linewidth}{@{\extracolsep{\fill}}llrrrrr@{}}
\toprule
Arm & Window & Lookup-69 & Lookup-30 & Lookup-SG & Physical 1 & Physical 2\\
\midrule
% Generated from independently recomputed E4 records.
H & R & 1.8007 & 1.0215 & 1.0186 & 0.9799 & 1.1618 \\
H & S & 2.1624$^{\dagger}$ & 1.3655$^{\dagger}$ & 1.3539$^{\dagger}$ & 1.3107$^{\dagger}$ & 1.2926$^{\dagger}$ \\
H+guard & R & 1.8007 & 1.0215 & 1.0186 & 1.0266 & 0.9678 \\
H+guard & S & 2.3375$^{\dagger}$ & 1.4256 & 1.4139 & 1.3930 & 1.3470 \\
C-1 & R & 2.1451 & 2.1696 & 2.1674 & 2.1119 & 2.1371 \\
C-1 & S & 2.2972 & 2.2966 & 2.2908 & 2.2300 & 2.2390 \\
C-2 & R & 2.6994$^{\dagger}$ & 1.4054 & 1.4025 & 1.3410 & 1.3429 \\
C-2 & S & 3.0270$^{\dagger}$ & 1.6321 & 1.6211 & 1.5520 & 1.5455 \\
N2 & R & 1.9102 & 1.8709 & 1.8704 & 1.8397 & 1.8390 \\
N2 & S & 2.0171 & 1.9517 & 1.9487 & 1.9248 & 1.9222 \\
N4 & R & 1.9650 & 1.9621 & 1.9621 & 1.9558 & 1.9551 \\
N4 & S & 2.0689$^{\dagger}$ & 2.0180 & 2.0180 & 2.0202 & 2.0191 \\
\bottomrule

\end{tabular*}
\end{table*}

\begin{table*}[t]
\centering
\small
\caption{Engine study physical resource use and offline on-time counts, shown for both
repeats. All online arrivals are on time; each offline denominator is 120.
H's two misses are unfinished jobs at cutoff. Occupancy covers the registered
window, including startup, draining, and stopping until verified release.}
\label{tab:e4resources}
\setlength{\tabcolsep}{4pt}
\par\vspace{4pt}
\begin{tabular*}{0.96\linewidth}{@{\extracolsep{\fill}}llrrrr@{}}
\toprule
Arm & Window & GPU seconds 1 & GPU seconds 2 & Offline 1 & Offline 2\\
\midrule
% Generated from independently recomputed E4 records.
H & R & 814.275 & 990.208 & 120/120 & 120/120 \\
H & S & 951.587 & 937.009 & 119/120 & 119/120 \\
H+guard & R & 836.218 & 802.155 & 120/120 & 120/120 \\
H+guard & S & 1041.078 & 1003.729 & 120/120 & 120/120 \\
PPO-1 & R & 2004.565 & 2026.170 & 120/120 & 120/120 \\
PPO-1 & S & 2015.445 & 2016.216 & 120/120 & 120/120 \\
PPO-2 & R & 1149.451 & 1150.762 & 120/120 & 120/120 \\
PPO-2 & S & 1129.373 & 1133.962 & 120/120 & 120/120 \\
PPO-7 & R & 1757.456 & 1759.098 & 120/120 & 120/120 \\
PPO-7 & S & 1765.317 & 1755.676 & 120/120 & 120/120 \\
PPO-9 & R & 1951.793 & 1951.098 & 120/120 & 120/120 \\
PPO-9 & S & 1982.853 & 1982.961 & 120/120 & 120/120 \\
\bottomrule

\end{tabular*}
\end{table*}

Lookup-69 falsely rejects eight feasible runs: both steady H+guard runs, all four
PPO-2 runs, and both steady PPO-9 runs. At arm level, requiring P1+ in all four
executions, Lookup-69 accepts two feasible arms and rejects four, three of which
are feasible. Lookup-30 and Lookup-SG each accept the five feasible arms and
reject H. Thus local startup calibration and Lookup-30 give identical feasibility
judgments here. Their total-cost MAEs are 0.0427 and 0.0468, respectively,
versus 0.5420 for Lookup-69. Excluding deadline penalties from both prediction and
observation gives operating-cost MAEs of 0.0427, 0.0468, and 0.3411 in
the same order. The cost-error difference is therefore not solely a
penalty effect.

H+guard passes all four runs and averages 920.795 occupied GPU seconds,
78.08\% below the all2 anchors; individual-run reductions span
75.21--80.90\%. These are window resource savings under the synthetic-API
contract. H's low occupancy alone would hide its two deadline failures.

Across valid formal runs, 255 of 256 ordinary in-window launches reach
readiness. Their median is 28.052\,s, p95 43.094\,s, and range
26.112--94.872\,s. The remaining launch is stopped by H+guard after
24.002\,s before readiness; its occupancy is charged, and it is excluded
from completed-startup statistics. Initial setup and invalid attempts
are excluded from this distribution. The calibration captures the typical
startup scale, while the physical record preserves its variability.

\paragraph{Retrospective startup sensitivity.}
We vary only Lookup-SG's startup proxy over 18 values from 20 to
69.219\,s, keeping policies and all other replay parameters fixed.
Against the same 24 physical outcomes, all nine sweep values from
20 through 41\,s retain 24/24 agreement; the 42\,s test gives 22/24
(Table~\ref{tab:e4sensitivity}). At Lookup-69's original startup value, replay
reproduces its eight false rejections. Exact agreement is therefore
not specific to the measured median or the 30\,s setting.
The source bundle includes the full grid and per-run comparisons.

\begin{table}[t]
\centering
\small
\caption{Selected values from the retrospective CPU sweep, plus the
measured-median reference (marked $*$). Startup values are rounded for
display. Each row reuses the same 24 engine policy runs; the sweep tests
discrete values, not a continuous tolerance interval.}
\label{tab:e4sensitivity}
\setlength{\tabcolsep}{3pt}
\par\vspace{4pt}
\begin{tabular*}{\linewidth}{@{\extracolsep{\fill}}rrrr@{}}
\toprule
Startup (s) & Agreement & \shortstack{False\\accepts} & \shortstack{False\\rejects}\\
\midrule
% Generated from the independently reproduced retrospective CPU sweep.
20 & 24/24 & 0/2 & 0/22 \\
27.568$^*$ & 24/24 & 0/2 & 0/22 \\
30 & 24/24 & 0/2 & 0/22 \\
41 & 24/24 & 0/2 & 0/22 \\
42 & 22/24 & 0/2 & 2/22 \\
69.219 & 16/24 & 0/2 & 8/22 \\
\bottomrule

\end{tabular*}
\end{table}

\paragraph{Technical attempts and verification.}
There are 28 attempts for 26 cells. PPO-1's two recovery executions each
encounter a control-loop timeout on their first attempt and are repeated
once with unchanged configurations under recorded authorization. The
historical causes remain unresolved; the failed attempts are retained.
An additional environment-retry provision was recorded during execution
but never used. H's two deadline-violating runs are not retried.
The CPU receiver verifies all file manifests and recomputes request
outcomes and resource/cost ledgers using the original frozen analyzer;
all 26 ledgers agree with the GPU return. The supplement's
\path{E4_RESULT_SPECIFICATION.json} records the source bindings, parameters,
and result definitions.
